\begin{figure*}[t]
  \scalebox{1.0} {
    \input{img/fig-series}
  }
  \subcaption{Using negated input $x_a x_b \bar{y}$}
  \sublabel{fig-series}
  \scalebox{1.0} {
    \input{img/fig-series2}
  }
  \subcaption{Using Boolean Fourier transform calculated $x_a x_b \bar{y}$}
  \sublabel{fig-series2}
  \caption{Phase oracle for $x_a x_b y \oplus x_a x_b \bar{y}$}
  \label{fig-series-all}
\end{figure*}

\section{Motivation}
\label{Mot}
\subsection{CNOT+T Circuit Optimization}
\label{Mot:CnotOpt}
CNOT+T circuits are notable because it is straightforward to optimize the T-count.
Note that since the phase is a scalar, the actions of phase gates on different
qubits all multiply together into the same scalar (essentially adding the powers
of $e^{i \pi}$). Therefore, the result of all phase gates in a circuit
can be computed as follows:

\begin{itemize}
\item For every phase gate, calculate the state driving it as a function of the
  input states, which will be a linear function of the same form as the
  $\chi_{\mathbf{k}}$ from the previous section.
\item Multiply the phase driven by the phase gate by this function.
\item Add the phases from all such similar functions.
\item Resynthesize the circuit from the resulting phase polynomial using
  the methods from Sec.~\ref{Pre}.
\end{itemize}

The result (sometimes called a circuit's {\it sum of paths}~\cite{bib-amy-cnot})
is that any phase gates driven by the same function will merge
together in the phase polynomial. This means that T-gates
can combine into an S-gate or Z-gate, reducing the total T-count.

This is important because the matroid partitioning algorithm from~\cite{bib-amy-matroid}
that was used to sequence Fig.~\ref{fig-toff-mark-matroid} uses this sum of paths
calculation as a preprocessing step. 

\begin{example}
  \label{ex-series}
  We attempt to calculate the phase polynomial for Fig.~\ref{fig-series}. This circuit
  is the phase oracle for $x_a x_b y \oplus x_a x_b \bar{y}$ implemented using
  the method detailed in Sec~\ref{Pre:OracleEsop}. It inserts
  Fig.~\ref{fig-toff-mark-matroid} to implement $x_a x_b y$. The second term has a negated
  input, so it inserts X on the $y$ qubit after the first circuit to negate $y$, again
  insert Fig.~\ref{fig-toff-mark-matroid} to implement the phase, and finally place another
  X on $y$ to return its value to the original.

  Since the first half of the circuit just implements $x_a x_b y$, we already know
  the phase polynomial for this from Eq.~\ref{eq-toff-bool}. The second term, we
  substitute $\bar{y}$ for $y$ in Eq.~\ref{eq-toff-bool}. We add the two phase
  polynomials together and get

  \begin{equation}
    \label{eq-series}
    \begin{aligned}
    % (actual polynomial here)
    \end{aligned}
  \end{equation}
\end{example}
